% META: {"id": "1SPE-SUITES-CO-035", "chapitre": "1SPE-SUITES", "type_objet": "corrige", "exercice_ref": "1SPE-SUITES-EX-035", "capacites": ["1SPE-SUITES-C2", "1SPE-SUITES-C6", "1SPE-SUITES-C7"], "capacites_codes": ["C2", "C6", "C7"], "sources_inspiration": [], "status": "generated"}
% BEGIN-VERIFY
% from sympy import symbols
% n = symbols('n', integer=True, nonneg=True)
% N0 = 1200
% r = 50
% N = N0 + r * n
% assert N.subs(n, 0) == 1200
% assert N.subs(n, 1) == 1250
% diff = N.subs(n, n+1) - N
% assert diff == 50
% # Population en 2031 : n = 11
% assert int(N.subs(n, 11)) == 1750
% # Somme S = N0 + N1 + ... + N11 : 12 termes, premier=1200, dernier=1750
% S = 12 * (1200 + 1750) // 2
% assert S == 17700
% # Seuil strict N_n > 2000 : 1200 + 50n > 2000 => n > 16
% assert int(N.subs(n, 16)) == 2000
% assert int(N.subs(n, 17)) == 2050
% END-VERIFY
\begin{corrige}{1SPE-SUITES-EX-035}

\textbf{Question 1 — Modélisation et nature arithmétique.}

La population augmente de $50$ habitants par an, donc :
\[
  N_{n+1} = N_n + 50, \quad N_0 = 1\,200.
\]

Pour tout $n \in \mathbb{N}$, $N_{n+1} - N_n = 50$.

\commentaireMarge{Différence constante $r=50$ : suite arithmétique.}

$(N_n)$ est une suite arithmétique de raison $r = 50$ et de premier terme $N_0 = 1\,200$, d'où :
\[
  N_n = 1\,200 + 50n.
\]

\medskip
\textbf{Question 2 — Population au début de l'année 2031.}

L'année 2031 correspond à $n = 2031 - 2020 = 11$.
\[
  N_{11} = 1\,200 + 50 \times 11 = 1\,200 + 550 = 1\,750 \text{ habitants}.
\]

\medskip
\textbf{Question 3 — Somme des douze relevés de population.}

La somme porte sur $12$ termes (de rang $0$ à $11$), de premier terme $N_0 = 1\,200$ et de dernier terme $N_{11} = 1\,750$.

\commentaireMarge{Formule indexée : $S = u_0 + \cdots + u_n = (n+1) \times \frac{u_0+u_n}{2}$.}

\[
  S = \frac{12 \times (N_0 + N_{11})}{2} = \frac{12 \times (1\,200 + 1\,750)}{2} = \frac{12 \times 2\,950}{2} = 6 \times 2\,950 = 17\,700.
\]

La somme des douze relevés vaut $17\,700$. Elle ne représente ni la population présente à une date donnée ni un nombre d'habitants distincts : une même personne peut être comptée dans plusieurs relevés annuels.

\medskip
\textbf{Question 4 — Analyse du programme Python.}

\textit{a) Explication ligne par ligne.}

\begin{itemize}
  \item \texttt{N = 1200} : initialisation de la population à $1\,200$ (valeur de $N_0$).
  \item \texttt{n = 0} : initialisation du compteur d'années à $0$.
  \item \texttt{while N <= 2000:} : on répète la boucle tant que la population n'a pas encore strictement dépassé $2\,000$.
  \item \texttt{N = N + 50} : on augmente la population de $50$ (passage à l'année suivante).
  \item \texttt{n = n + 1} : on incrémente le compteur d'années.
  \item \texttt{print("Annee :", 2020 + n)} : on affiche l'année calendaire.
  \item \texttt{print("Population :", N)} : on affiche la population atteinte.
\end{itemize}

\textit{b) Résultat sans exécution.}

On cherche le plus petit entier $n$ tel que $N_n > 2\,000$ :
\[
  1\,200 + 50n > 2\,000 \iff 50n > 800 \iff n > 16.
\]

\commentaireMarge{Arrêt à $n=17$, affiche $2037$.}

Le plus petit entier $n$ vérifiant cette condition est $n = 17$.

Le programme affiche \texttt{Annee : 2037} puis \texttt{Population : 2050}.

\textit{c) Vérification.}
\[
  N_{16} = 1\,200 + 50 \times 16 = 2\,000 \quad \text{(la boucle continue encore)}.
\]
\[
  N_{17} = 1\,200 + 50 \times 17 = 2\,050 > 2\,000 \quad \text{(la boucle s'arrête)}. \checkmark
\]

La ville dépassera $2\,000$ habitants au début de l'année 2037.

\baremeIndicatif{Q1 : 2 pts (modeliser, raisonner : relation de récurrence 1 pt + terme général 1 pt) ; Q2 : 1 pt (calculer : $N_{11}$ avec identification du rang) ; Q3 : 2 pts (calculer : somme 1 pt + interprétation 1 pt) ; Q4a : 2 pts (programmer, communiquer : explication ligne par ligne) ; Q4b : 2 pts (raisonner, calculer : résolution algébrique et résultat) ; Q4c : 1 pt (calculer : vérification $N_{16}$ et $N_{17}$)}
\end{corrige}
